\documentclass[11pt,a4paper]{article}
\usepackage[margin=2.2cm]{geometry}
\usepackage{booktabs}
\usepackage{graphicx}
\usepackage{hyperref}
\usepackage{xcolor}
\usepackage{amsmath}
\usepackage{multirow}
\usepackage{tabularx}
\usepackage[T1]{fontenc}

\setlength{\parskip}{0.4em}
\setlength{\parindent}{0em}

\title{\vspace{-1.5cm}\textbf{Assignment~2: Learning from Click-Logs}\\
\large CS4.406 --- Information Retrieval \& Extraction \\
\normalsize GBDT Re-Ranker Track}
\author{Shubham Paliwal}
\date{September 2026}

\begin{document}
\maketitle
\vspace{-0.6cm}

%=========================================================================
\section{System Overview}
%=========================================================================

We build a two-stage \emph{retrieve-then-rank} pipeline for news recommendation on
\textbf{EB-NeRD} (Danish, Ekstra Bladet) and \textbf{MIND} (English, Microsoft News),
extending the Assignment~1 codebase (\texttt{IRE-Assignment1}) rather than rewriting it.
All new logic lives in \texttt{src/ire\_a2/} and \texttt{scripts/}, reusing A1's
from-scratch BM25 index, embedding index, and evaluation harness directly.

\textbf{Stage 1: Candidate Retrieval.} Three baselines score each impression's served
candidates: \emph{B1 (Lexical)} is A1's dict-based inverted-index BM25 over
title+abstract. \emph{B2 (Semantic)} is cosine similarity between a mean-pooled user
history embedding and each candidate embedding (sentence-transformer model
\emph{paraphrase-multilingual-MiniLM-L12-v2}).
\emph{B3 (Learned Hybrid)} fuses B1 and B2 with a weight $\alpha$ grid-searched over
$[0,1]$ (step $0.05$) to maximise \emph{mean per-impression} AUC on a train-split
sample:
\[
  s_{\text{hybrid}} = \alpha \cdot s_{\text{BM25}} + (1-\alpha) \cdot s_{\text{embed}}
\]
On both datasets the learned weight collapsed to $\alpha=0.00$ (pure embeddings) --
plain embeddings already maximised the fit-sample objective, so B3 is numerically
identical to B2 here (a real, not fabricated, degenerate result -- unlike the NRMS
track's B3, which found a non-degenerate blend on this data).

\textbf{Stage 2: GBDT Re-Ranker.} A LightGBM LambdaRank model scores each candidate
from 13 hand-crafted features (Section~2). Grouped by impression, trained with fixed
hyperparameters for local evaluation (Sections 3--6) and separately hyperparameter-tuned
for the large-scale Codabench submission (Section~7).

\textbf{Stack.} Polars for all tabular processing, a from-scratch dict-based inverted
index for BM25 (not a library), \texttt{sentence-transformers} for embeddings,
LightGBM for the re-ranker. Scale: Sections 3--6 use \texttt{demo} (EB-NeRD, 11{,}777
articles) / \texttt{small} (MIND, 65{,}238 articles) for fast, reproducible local
iteration; Section~7 uses the real \texttt{large} bundles (125{,}541 / 130{,}379
articles) for the actual Codabench submission.

%=========================================================================
\section{Feature Engineering (Q1)}
%=========================================================================

Thirteen features are computed per (user, candidate) pair by \texttt{FeatureBuilder}
(\texttt{src/ire\_a2/}\allowbreak\texttt{features.py}), grouped into four families:

\begin{table}[h]
\centering\small
\begin{tabularx}{\textwidth}{lXl}
\toprule
\textbf{Family} & \textbf{Feature} & \textbf{Source} \\
\midrule
\multirow{2}{*}{Retrieval}
 & \texttt{bm25\_score} --- BM25 relevance of candidate to user's history titles & B1 \\
 & \texttt{embedding\_score} --- cosine similarity, mean-pooled history vs. candidate & B2 \\
\midrule
Position
 & \texttt{candidate\_position} --- candidate's 0-indexed position in the served list & impression \\
\midrule
\multirow{4}{*}{Click-history}
 & \texttt{hist\_click\_count} --- clicks strictly before this impression (capped at 20) & history \\
 & \texttt{hist\_recency\_weighted\_count} --- $\sum e^{-\ln2 \cdot \Delta t / 72h}$ over history & history \\
 & \texttt{hist\_category\_match} / \texttt{\_weighted} --- (recency-weighted) count of history & history \\
 & \quad clicks sharing the candidate's category & \\
\midrule
\multirow{3}{*}{Article}
 & \texttt{candidate\_popularity} --- $\log(1{+}\text{train-split click count})$ & articles \\
 & \texttt{freshness\_hours} / \texttt{has\_freshness} --- hours since publish (EB-NeRD only) & articles \\
 & \texttt{hist\_avg\_dwell\_seconds}, \texttt{session\_ordinal}, & articles \\
 & \quad \texttt{session\_click\_count\_so\_far} --- EB-NeRD-only session/dwell signal & \\
\bottomrule
\end{tabularx}
\caption{Feature set (13 features). \texttt{candidate\_popularity} is frozen to
train-split-only counts to prevent leakage (Section~6). MIND provides no session id or
per-click dwell time at all -- a real data limitation, not a modelling choice -- so the
three session/dwell features are always 0 for MIND.}
\label{tab:features}
\end{table}

Every feature depends only on \texttt{UserHistoryIndex.recent(user\_id, cutoff\_ts)}
(strictly-before-cutoff click history) or a train-split-only frozen popularity dict --
never on validation/test data. This is the behaviour-window boundary Section~6 tests
directly.

%=========================================================================
\section{Re-Ranker Architecture (Q2)}
%=========================================================================

\textbf{Model.} LightGBM's \texttt{LGBMRanker} with the \texttt{lambdarank} objective,
grouped by impression\_id (LightGBM's native pairwise-preference ranking loss, not a
from-scratch implementation -- unlike BM25, gradient-boosted trees are not something
this project reimplements). Impressions with an all-clicked or all-unclicked label set
are dropped before training (LambdaRank needs both classes to form a pair, the same
rule used for AUC/nDCG).

\textbf{Hyperparameters.} Sections 3--6 use fixed, untuned defaults
($\texttt{num\_leaves}{=}31$, $\texttt{learning\_rate}{=}0.05$,
$\texttt{n\_estimators}{=}200$, $\texttt{min\_child\_samples}{=}5$) deliberately, to
keep the baseline-vs-re-ranker comparison simple and reproducible. Section~7 describes
a real hyperparameter search used only for the large-scale Codabench submission.

%=========================================================================
\section{Experimental Results (Q2/Q3)}
%=========================================================================

\subsection{Baseline \& Re-Ranker Performance}

Evaluated on a 5{,}000-impression sample of each dataset's validation split (seed 42,
matching Assignment~1's own evaluation convention). 95\% bootstrap confidence
intervals ($n=1000$) reported.

\begin{table}[h]
\centering\small
\begin{tabular}{ll cccc}
\toprule
\textbf{Dataset} & \textbf{Method} & \textbf{AUC} & \textbf{MRR} & \textbf{nDCG@5} & \textbf{nDCG@10} \\
\midrule
\multirow{4}{*}{\rotatebox{90}{\textbf{EB-NeRD}}}
 & B1 BM25       & 0.494 & 0.312 & 0.341 & 0.427 \\
 & B2 Embedding  & 0.539 & 0.338 & 0.375 & 0.454 \\
 & B3 Hybrid ($\alpha{=}0.00$) & 0.539 & 0.338 & 0.375 & 0.454 \\
 & \textbf{GBDT} & \textbf{0.559} & 0.324 & 0.372 & \textbf{0.455} \\
\midrule
\multirow{4}{*}{\rotatebox{90}{\textbf{MIND}}}
 & B1 BM25       & 0.553 & 0.293 & 0.269 & 0.331 \\
 & B2 Embedding  & \textbf{0.633} & \textbf{0.340} & \textbf{0.325} & \textbf{0.384} \\
 & B3 Hybrid ($\alpha{=}0.00$) & 0.633 & 0.340 & 0.325 & 0.384 \\
 & GBDT          & 0.594 & 0.340 & 0.318 & 0.373 \\
\bottomrule
\end{tabular}
\caption{Validation-sample results (demo/small scale). Bold marks the best method per
dataset per metric. On EB-NeRD GBDT wins every metric; on MIND, B2/B3 (embeddings)
wins every metric -- reported plainly, not the outcome hoped for.}
\label{tab:results}
\end{table}

\subsection{Paired Bootstrap Ablation (Q3)}

\begin{table}[h]
\centering\small
\begin{tabular}{ll cc l}
\toprule
\textbf{Dataset} & \textbf{Comparison} & \textbf{$\Delta$AUC (GBDT $-$ baseline)} & \textbf{95\% CI} & \textbf{Sig.?} \\
\midrule
\multirow{3}{*}{EB-NeRD}
 & GBDT vs.\ B1 BM25   & $+0.065$ & $[+0.053, +0.076]$ & YES \\
 & GBDT vs.\ B2 Embed  & $+0.020$ & $[+0.008, +0.031]$ & YES \\
 & GBDT vs.\ B3 Hybrid & $+0.020$ & $[+0.008, +0.031]$ & YES \\
\midrule
\multirow{3}{*}{MIND}
 & GBDT vs.\ B1 BM25   & $+0.041$ & $[+0.029, +0.052]$ & YES \\
 & GBDT vs.\ B2 Embed  & $-0.040$ & $[-0.049, -0.031]$ & YES (loss) \\
 & GBDT vs.\ B3 Hybrid & $-0.040$ & $[-0.049, -0.031]$ & YES (loss) \\
\bottomrule
\end{tabular}
\caption{Paired bootstrap ablation, AUC. GBDT beats BM25 significantly and
consistently on both datasets. On EB-NeRD it also beats embeddings significantly on
AUC (MRR loses, $-0.014$, $[-0.023,-0.004]$, significant -- also reported, not
cherry-picked). On MIND it loses to embeddings significantly on AUC and is not
significantly different on MRR ($-0.001$, $[-0.009,+0.009]$).}
\label{tab:ablation}
\end{table}

The honest read: GBDT's clearest, most consistent win is over BM25 -- a real,
significant, non-trivial improvement on both datasets. It does \emph{not} cleanly beat
a strong embedding baseline; the direction of the effect even flips between datasets
(a real gain on EB-NeRD's AUC, a real loss on MIND's AUC). Section~9 (Comparative
Study) discusses this pattern's persistence at Codabench's real test scale.

%=========================================================================
\section{Slicing \& Beyond-Accuracy (Q5)}
%=========================================================================

Cold-start threshold fixed at \texttt{history\_len}$<5$; head/tail split on the
clicked article's train-split popularity (75th percentile threshold).

\begin{table}[h]
\centering\small
\begin{tabular}{l cccc ccc}
\toprule
& \multicolumn{4}{c}{\textbf{AUC by Slice (GBDT)}} & \multicolumn{3}{c}{\textbf{Beyond-Accuracy (overall)}} \\
\cmidrule(lr){2-5}\cmidrule(lr){6-8}
\textbf{Dataset} & Cold & Warm & Tail & Head & Diversity@5 & Novelty@5 & Coverage@5 \\
\midrule
EB-NeRD & --- & 0.559 & 0.554 & 0.876 & 0.844 & 14.39 & 0.137 \\
MIND    & 0.593 & 0.594 & 0.535 & 0.739 & 0.894 & 14.94 & 0.012 \\
\bottomrule
\end{tabular}
\caption{EB-NeRD's demo-scale val sample contained no cold-start impressions at this
threshold (n=0), reported as such rather than silently omitted. Both datasets show a
large, consistent head/tail gap -- GBDT is substantially better at ranking popular
articles correctly, which \texttt{candidate\_popularity} (the top-gain feature in every
run) plausibly explains directly.}
\label{tab:slicing}
\end{table}

MIND's coverage is far lower than EB-NeRD's (1.2\% vs.\ 13.7\%) -- consistent with
MIND's much larger catalogue (65K vs.\ 11.8K articles at this scale) relative to what a
top-5 ranked list can ever surface.

%=========================================================================
\section{Serving \& Scale Analysis (Q4)}
%=========================================================================

\begin{table}[h]
\centering\small
\begin{tabular}{lcc}
\toprule
\textbf{Component} & \textbf{EB-NeRD (demo)} & \textbf{MIND (small)} \\
\midrule
BM25 index & 5.9\,MB & 42.5\,MB \\
Embedding matrix & 18.1\,MB & 100.2\,MB \\
Feature store (popularity+category dicts) & 0.25\,MB & 1.24\,MB \\
\midrule
Candidate retrieval (embedding search, $K{=}200$) & 0.61\,ms & 2.92\,ms \\
Feature build & 1.77\,ms & 2.64\,ms \\
GBDT scoring & 0.93\,ms & 1.05\,ms \\
\textbf{Total p99 latency} & \textbf{3.64\,ms} & \textbf{7.74\,ms} \\
Single-core capacity & 274\,req/s & 129\,req/s \\
Cost / 1{,}000 queries (\$0.05/vCPU-hr) & \$0.00005 & \$0.00011 \\
\bottomrule
\end{tabular}
\caption{Measured single-request latency, full pipeline (candidate generation through
GBDT score), 200 repeated runs per dataset.}
\label{tab:serving}
\end{table}

\textbf{10$\times$ scale bottleneck.} Candidate retrieval is brute-force cosine
similarity over the full catalogue -- $O(\text{catalogue size})$ per query. At 10$\times$
scale this is the first thing to break, well before feature building or GBDT scoring
(the cheapest stage today, and inherently the easiest to keep cheap since LightGBM's
\texttt{predict()} is already vectorised over the whole candidate batch). The fix is an
ANN index (FAISS/ScaNN), exactly as A1's own embedding-search code already flags. Feature
building (a per-candidate Python loop) is the next bottleneck once retrieval is
ANN-accelerated; memory (tens to low hundreds of MB even at this scale) is not a
near-term concern.

%=========================================================================
\section{Anti-Leakage Measures (Q9)}
%=========================================================================

\textbf{Behaviour-window boundary test.} A dedicated unit test
(\texttt{tests/}\allowbreak\texttt{test\_a2\_features.py}) constructs a toy user history containing a click
exactly at the query cutoff timestamp and asserts \texttt{UserHistoryIndex.recent()}
excludes it -- re-verified live in the companion notebook, not just referenced.

\textbf{With/without-leak ablation.} \texttt{candidate\_popularity} is deliberately
recomputed using train\,+\,validation hindsight click counts (via a
\texttt{popularity\_override} hook on \texttt{FeatureBuilder}) instead of the safe
train-only default, and a second GBDT is trained and compared via paired bootstrap.

\begin{table}[h]
\centering\small
\begin{tabular}{l cccc}
\toprule
\textbf{Dataset} & \textbf{Metric} & \textbf{$\Delta$ (with $-$ without leak)} & \textbf{95\% CI} & \textbf{Sig.?} \\
\midrule
\multirow{4}{*}{EB-NeRD}
 & AUC     & $+0.149$ & $[+0.139, +0.159]$ & YES \\
 & MRR     & $+0.140$ & $[+0.130, +0.151]$ & YES \\
 & nDCG@5  & $+0.155$ & $[+0.144, +0.166]$ & YES \\
 & nDCG@10 & $+0.123$ & $[+0.115, +0.132]$ & YES \\
\midrule
\multirow{4}{*}{MIND}
 & AUC     & $+0.068$ & $[+0.058, +0.077]$ & YES \\
 & MRR     & $+0.048$ & $[+0.038, +0.057]$ & YES \\
 & nDCG@5  & $+0.047$ & $[+0.038, +0.056]$ & YES \\
 & nDCG@10 & $+0.049$ & $[+0.041, +0.057]$ & YES \\
\bottomrule
\end{tabular}
\caption{The leak inflates every metric, on both datasets, with a large and clearly
significant effect -- a clean demonstration of what a real leak looks like (contrast
with a near-zero or ambiguous effect, which would be much harder to interpret as a real
leak vs. noise).}
\label{tab:q9}
\end{table}

%=========================================================================
\section{Hyperparameter Tuning (Q6, large-scale submission only)}
%=========================================================================

Sections 3--6 deliberately use fixed hyperparameters. For the actual Codabench
submission, a random search (12 trials, not an exhaustive $4^4{=}256$-combination
grid) over \texttt{num\_leaves}, \texttt{learning\_rate}, \texttt{n\_estimators},
\texttt{min\_child\_samples} was run against each dataset's \emph{real} validation
split (not the tuning-train data), scored by AUC.

\begin{table}[h]
\centering\small
\begin{tabularx}{\textwidth}{l X r r}
\toprule
\textbf{Dataset} & \textbf{Winning config} & \textbf{AUC} & \textbf{Search time} \\
\midrule
MIND & \texttt{num\_leaves=15, lr=0.02, n\_est=200, min\_child=5} & 0.640 & 3.5\,min \\
EB-NeRD & \texttt{num\_leaves=15, lr=0.02, n\_est=300, min\_child=3} & 0.699 & 2.0\,min \\
\bottomrule
\end{tabularx}
\caption{Both winning configurations use simpler trees (\texttt{num\_leaves=15}) and a
lower learning rate than the fixed default -- a consistent pattern across both
datasets. The search was cheap enough (under 4 minutes each) that both datasets got
tuned, not just the one originally requested.}
\label{tab:tuning}
\end{table}

The winning configuration was retrained on a full-scale (50{,}000-impression) training
sample and used for the actual submission (Section~8) -- not the small tuning-train
sample used only to compare candidate configurations cheaply.

%=========================================================================
\section{Comparative Study: GBDT vs.\ NRMS}
%=========================================================================

A teammate independently built an NRMS re-ranker track (Option B) using a
from-scratch reimplementation of the A1 pipeline with different libraries
(\texttt{rank\_bm25}, NLTK, different embedding models) and a different evaluation
sample size (full validation split vs.\ this track's 5{,}000-impression sample) --
\textbf{cross-track raw numbers are not a controlled comparison of GBDT vs.\ NRMS
alone}. Full detail in \texttt{A2\_TRACK\_COMPARISON.md}.

\begin{table}[h]
\centering\small
\begin{tabular}{ll cccc}
\toprule
\textbf{Dataset} & \textbf{Re-Ranker} & \textbf{AUC} & \textbf{MRR} & \textbf{nDCG@5} & \textbf{nDCG@10} \\
\midrule
\multirow{2}{*}{EB-NeRD}
 & GBDT  & 0.559 & 0.324 & 0.372 & 0.455 \\
 & NRMS  & \textbf{0.563} & \textbf{0.349} & \textbf{0.390} & \textbf{0.466} \\
\midrule
\multirow{2}{*}{MIND}
 & GBDT  & 0.594 & 0.340 & 0.318 & 0.373 \\
 & NRMS  & \textbf{0.656} & \textbf{0.349} & \textbf{0.333} & \textbf{0.396} \\
\bottomrule
\end{tabular}
\caption{Cross-track comparison (see caveats above). Within-track comparisons are more
reliable: NRMS shows a cleaner, more consistent win over its own baselines than GBDT
does over embeddings.}
\label{tab:cross}
\end{table}

The GBDT track additionally contributed a \textbf{quantified with/without-leak
ablation} for Q9 (Table~\ref{tab:q9}); the NRMS track identified the same class of risk
but did not measure it with a paired before/after comparison.

%=========================================================================
\section{Codabench Submissions (Q6)}
%=========================================================================

\textbf{Large-scale retraining.} The submission model is \emph{not} the demo/small
model from Sections 3--6 -- a fresh GBDT was trained on a 50{,}000-impression sample of
each dataset's real \texttt{large} train split, using the tuned hyperparameters from
Section~7, so that \texttt{candidate\_popularity} (a raw click count) reflects the same
scale the real unlabeled test set was drawn from.

\textbf{Two real engineering problems, found and fixed.} EB-NeRD's large history data
OOM'd this project's development machine (16GB RAM) via a global explode+dedup already
present in Assignment~1's shared \texttt{clean\_user\_history()} -- fixed at the source
(verified numerically identical to the old implementation on real data first), then a
second fix limited history construction to only the users actually needed by the
training sample rather than the full population. Separately, 200{,}000 real EB-NeRD
test rows (a deliberate, round number) share \texttt{raw\_impression\_id == "0"};
naive grouping by that id silently collapses all of them into a single output line --
verified directly with a toy example before fixing it (a synthetic per-row grouping key,
with the original id reattached only for the line actually written).

\begin{table}[h]
\centering\small
\begin{tabularx}{\textwidth}{l r r X}
\toprule
\textbf{Dataset} & \textbf{Written} & \textbf{Expected} & \textbf{Match?} \\
\midrule
MIND    & 2{,}370{,}727  & 2{,}370{,}727  & Exact \\
EB-NeRD & 13{,}536{,}710 & 13{,}536{,}710 & Exact, including the 200{,}000 id-0 rows, verified as separate lines \\
\bottomrule
\end{tabularx}
\caption{Row counts verified against A2.pdf's stated test-set sizes before upload.}
\end{table}

\textbf{Real leaderboard results.}

\begin{table}[h]
\centering\small
\begin{tabularx}{\textwidth}{l X r}
\toprule
\textbf{Dataset} & \textbf{Submission} & \textbf{Score} \\
\midrule
\multirow{2}{*}{MIND}
 & \texttt{mind\_prediction.zip} --- A1, embedding heuristic, no training & 0.6544 \\
 & \texttt{mind\_a2\_gbdt\_predictions.zip} --- this track, tuned GBDT & 0.5863 \\
\midrule
EB-NeRD & \texttt{ebnerd\_a2\_gbdt\_predictions.zip} & pending \\
\bottomrule
\end{tabularx}
\caption{MIND's real test-set score is reported plainly, including that it is
\emph{lower} than A1's untrained embedding heuristic -- discussed in Section~10, not
omitted. EB-NeRD's leaderboard outcome was not yet resolved as of writing this note;
the submission file itself is verified correct (row count, format, filename) per the
table above.}
\end{table}

%=========================================================================
\section*{Observations \& Discussion}
%=========================================================================

\textbf{GBDT's real, defensible win is over BM25}, not over embeddings. Both locally
(Table~\ref{tab:ablation}) and at Codabench's real test scale, GBDT does not cleanly
beat a strong embedding baseline on MIND -- the real leaderboard gap (0.5863 vs.\
0.6544) is if anything \emph{larger} than the local validation gap suggested
($-0.040$ AUC locally). The most likely explanation is training-data scale: GBDT here
trains on a 50{,}000-impression sample -- 2.2\% of MIND-large's 2.2M-impression train
split -- while A1's embedding heuristic needs no training data at all and is immune to
sample-size effects entirely. Hyperparameter tuning (Section~7) found a real,
cheaply-obtained local improvement (MIND AUC $0.624\to0.640$ on the tuning validation
sample) but did not close this gap at real test scale, suggesting the limitation is
sample size and/or feature design rather than tree configuration.

\textbf{The engineering cost of the large-scale run was substantial and real}, not
incidental -- two genuine bugs (one inherited from shared A1 infrastructure, one newly
discovered) were found only by actually running the full pipeline against real
production-scale data, not by local demo-scale testing. Both were root-caused and fixed
at the source rather than worked around, and both fixes were verified against real data
before being trusted.

\textbf{EB-NeRD's own local results are more encouraging for GBDT} than MIND's --
GBDT beats every baseline on AUC there, including embeddings, both locally and (subject
to the pending leaderboard result) presumably at scale too. The freshness/session
features available only on EB-NeRD (Section~2) may be part of why: EB-NeRD's feature
set has real signal MIND's simply does not have access to.

%=========================================================================
\section*{AI Usage Declaration}
%=========================================================================

Claude (Anthropic, Claude Code) was used for essentially all code generation,
architectural design, debugging, and drafting in this repository, including this
design note. A detailed, chronological \texttt{AI\_USAGE\_LOG.md} in the repository
root records specific prompts, decisions, and the student's direction and verification
at each stage -- including the two real bugs described in Section~8, the honest
reporting of GBDT's losses to embeddings (Sections 4, 10), and the real Codabench
scores (Section~8). Final review, all real experimental execution (including all
Codabench submissions), and interpretation of results were performed by the student.

\begin{thebibliography}{9}
\bibitem{ke2017lightgbm} G.~Ke, Q.~Meng, T.~Finley, T.~Wang, W.~Chen, W.~Ma, Q.~Ye, and
T.-Y.~Liu, ``LightGBM: A highly efficient gradient boosting decision tree,'' in
\emph{Advances in Neural Information Processing Systems}, 2017.
\bibitem{kruse2024ebnerd} J.~Kruse, K.~Lindskow, S.~Kalloori, M.~Polignano, C.~Pomo,
A.~Srivastava, A.~Uppal, M.~R.~Andersen, and J.~Frellsen, ``EB-NeRD: A large-scale
dataset for news recommendation,'' in \emph{Proc. ACM RecSys Challenge 2024}, 2024.
\bibitem{wu2020mind} F.~Wu, Y.~Qiao, J.-H.~Chen, C.~Wu, T.~Qi, J.~Lian, D.~Liu,
X.~Xie, J.~Gao, W.~Wu, and M.~Zhou, ``MIND: A large-scale dataset for news
recommendation,'' in \emph{Proc. ACL}, 2020, pp.~3597--3606.
\end{thebibliography}

\end{document}
